\subsection{Outlook}
\label{subsec:ausblick}

In the v1.0.0 candidate state, the historical Ruff, exception, and documentation
discrepancies are corrected: \path{tools/.venv} is profile-independent,
exception boundaries are enforced, and documentation synchronization is
checked. As the next release evidence, the master, all five profiles,
PostgreSQL, Rust/Tauri, and Windows must be checked on the same final commit.
Compose startup with health and readiness checks, Release
Validation, and Branch Protection should supplement the technical evidence
without claiming production readiness.

In the medium term, additional real negative Ruff, ESLint, and Rust fixtures
beyond the regression cases for v1.0.0 would be useful. Existing JSON reports
could be extended with SARIF, pull-request annotations, and trends. False
positives and the appropriateness of thresholds should be observed by profile
or project type. Other possible baseline extensions remain automated shared
contracts, stricter capability selection, and controlled synchronization of
generated projects.

Specific products still require signing, notarization, an updater, credentials,
deployment, authentication, and RBAC. In the long term, a comparative study
across projects should measure onboarding, time to the first build, setup
errors, refactoring frequency, defects, and maintenance effort. A controlled
comparison of agent-assisted development with and without the gate could test
whether technical rule detection has effects beyond the observed snapshot.
